\subsection{道路连续映射}

定义 5. --- 设 X 与 Y 为拓扑空间。映射 \(f: X \to Y\) 称为道路连续的，如果对 X 中的每条道路 \(c\)，映射 \(f \circ c: I \to Y\) 都连续。

注. --- 设空间 X 道路连通。设 x 为 X 的一点。为使 f 道路连续，只需对每条以 x 为起点的道路 c，映射 \(f \circ c\) 连续。事实上，设 d 为 X 中任意道路，c 为 X 中以 x 为起点、\(d(0)\) 为终点的道路。若映射 \(f \circ (c * d)\) 连续，则映射 \(f \circ d\) 也连续，因为有 \(f \circ d(t) = f \circ (c * d)((t + 1)/2)\)。

连续映射是道路连续的。反之不总成立；下面的命题 12 与 13 给出了断言道路连续映射连续的判据。

命题 12. --- 设 X 与 Y 为拓扑空间，\(f: \mathrm{X} \to \mathrm{Y}\) 为映射。假设空间 X 局部道路连通，且 X 的每一点都有可数的基本邻域系。若映射 \(f\) 是道路连续的，则它连续。

由（TG, IX, p. 18），只需证明：对 X 的每一点 \(x\) 与每个序列 \((x_n)_{n \geqslant 1}\)（由 X 中趋于 \(x\) 的点组成），序列 \((f(x_n))_{n \geqslant 1}\) 趋于 \(f(x)\)。删去序列 \((x_n)_{n \geqslant 1}\) 的开头若干项，可化归为序列的各项都属于 \(x\) 在 X 中的同一道路连通邻域的情形。由下面的引理，于是存在道路 \(c: \mathbf{I} \to \mathbf{X}\)，满足 \(c(0) = x\) 与 \(c(1/n) = x_n\)（对 \(n \geqslant 1\)）。若映射 \(f\) 道路连续，则映射 \(f \circ c\) 连续，而元素 \(f(x) = f(c(0))\) 是序列 \((f(c(1/n)))_{n \geqslant 1}\) 即序列 \((f(x_n))_{n \geqslant 1}\) 的极限，推论得证。

引理. --- 设 X 为连通且局部道路连通的拓扑空间，x 为 X 的一点。假设点 x 有可数的基本邻域系。那么，对每个趋于 x 的 X 的点的序列 \((x_{n})_{n\geqslant1}\)，存在 X 中的道路 c，满足 \(c(0)=x\) 与 \(c(1/n)=x_{n}\)（对 \(n\geqslant1\)）。

设 \((\mathrm{W}_{m})_{m\geqslant1}\) 为 x 的基本邻域系。令 \(V_{0}=X\)，且对每个 \(m\geqslant1\)，设 \(V_{m}\) 为包含于 \(V_{m-1}\cap W_{m}\) 中的 x 的道路连通邻域。

对每个整数 \(n \geqslant 1\)，用 \(m_n\) 表示最大整数 \(m \leqslant n\)，它满足 \(x_k \in V_m\)（对所有 \(k \geqslant n\)）。序列 \((m_n)_{n \geqslant 1}\) 按定义是递增的；由于序列 \((x_n)_{n \geqslant 1}\) 趋于 \(x\)，它趋于无穷。对每个整数 \(n \geqslant 1\)，设 \(c_n: \mathbf{I} \to V_{m_n}\) 为以 \(x_{n+1}\) 为起点、\(x_n\) 为终点的道路（在 \(V_{m_n}\) 中）。如下定义映射 \(c: \mathbf{I} \to X\)：令 \(c(0) = x\)，且 \(c(t) = c_n(n(n+1)t - n)\)（当 \(1/(n+1) < t \leqslant 1/n\) 时）。有 \(c(1/n) = x_n\)（对所有 \(n \geqslant 1\)）。因此映射 \(c\) 在每个形如 \([1/(n+1), 1/n]\)（\(n \geqslant 1\)）的区间上连续，从而在区间 \([0,1]\) 上连续。若 \(t \leqslant 1/n\)，则点 \(c(t)\) 属于 \(V_{m_n}\)；因此映射 \(c\) 在 0 处连续。

命题 13. --- 设 \(p \colon \mathrm{E} \to \mathrm{B}\) 为分离的平展态射，\(s \colon \mathrm{B} \to \mathrm{E}\) 为 \(p\) 的截面。若空间 B 局部道路连通且截面 \(s\) 道路连续，则它连续。

设 \(b\) 为 B 的一点；我们来证明 \(s\) 在点 \(b\) 处连续。由于 \(p\) 是平展映射，存在 \(b\) 的邻域 V 与定义在 V 中的连续局部截面 \(s'\)（\(p\) 的），满足 \(s'(b) = s(b)\)（I, p. 33, 命题 9）。由于 B 局部道路连通，可以假设 V 道路连通。对 V 中每条道路 \(c\)（以 \(b\) 为起点），映射 \(s \circ c\) 与 \(s' \circ c\) 是映射 \(c \colon \mathbf{I} \to \mathbf{B}\) 到 E 中的两个连续提升，且有 \(s \circ c(0) = s' \circ c(0) = s(b)\)。由 I, p. 34 的推论 1，有 \(s \circ c = s' \circ c\)，特别地有 \(s \circ c(1) = s' \circ c(1)\)。由于 V 的每一点都是 V 中从 \(b\) 出发的道路的终点，映射 \(s\) 与 \(s'\) 在 V 上一致。因此映射 \(s\) 在 V 中连续。

推论. --- 设 B 为拓扑空间，(E, p) 与 (E', p') 为两个 B-空间。假设映射 p 是平展且分离的，空间 E' 局部道路连通。那么每个道路连续映射 \(f: E' \to E\)（满足 \(p \circ f = p'\)）都是连续的。

映射 \(pr_{1}: E'\times_{B} E \rightarrow E'\) 是平展且分离的（I, p. 31, 命题 8 及 I, p. 27, 命题 4），而映射 \(x \mapsto (x, f(x))\) 是道路连续的截面。由命题 13，它连续，因此 f 连续。

